% =========================================================================
% Seguimiento de riesgos — 16-08-2026 (S4, cuarta retrospectiva).
% Primera revisión COMPLETA por cadencia: críticos (semanal) + significativos
% (cada 2: S2, S4) + moderados (cada 4: primera). Documento independiente de
% la memoria; creado copiando TEMPLATE.tex (ver README.md de esta carpeta).
% Compilar DESDE ESTA CARPETA (2 pases, sin biber):
%   pdflatex --enable-installer -interaction=nonstopmode -file-line-error
%     -output-directory=build 2026-08-16-S4.tex
% =========================================================================
\documentclass[11pt,a4paper]{article}

% =========================================================================
% Preámbulo COMPARTIDO de los documentos independientes de riesgos
% (docs/thesis/plan-riesgos/ y docs/thesis/seguimiento-riesgos/).
%
% Estos documentos NO forman parte de main.tex (la memoria): son entregables
% sueltos en clase article, con el diseño heredado de los .tex de Gestión de
% Proyectos del autor, pero adaptados a la configuración global de la
% memoria (Carlito, márgenes 2,5 cm, parskip 6 pt, títulos de la guía).
% Un solo preámbulo para toda la familia: un ajuste de estilo aquí se
% propaga a todos los seguimientos futuros (uno por sprint hasta S37).
%
% RUTAS: los \input son relativos a la carpeta del documento, no a esta;
% compilar SIEMPRE desde la carpeta del .tex (plan-riesgos/ o
% seguimiento-riesgos/):
%   pdflatex --enable-installer -interaction=nonstopmode -file-line-error
%     -output-directory=build <documento>.tex        (2 pases; sin biber)
% =========================================================================

% --- Codificación (idéntico a packages.tex de la memoria) ---
\usepackage[utf8]{inputenc}
\usepackage[T1]{fontenc}

% --- Idioma: solo español (aquí no hay resumen trilingüe) ---
% es-noshorthands por coherencia con la memoria: acentos en UTF-8, sin
% caracteres activos (el conflicto con biblatex-apa no aplica aquí, pero
% mantener una única configuración de babel evita sorpresas al copiar texto
% entre la memoria y estos documentos).
\usepackage[spanish, es-tabla, es-noshorthands]{babel}

% --- Márgenes: 2,5 cm laterales e inferior, como la memoria ---
% El margen superior se reparte para alojar la cabecera enmarcada del
% diseño de referencia (2 cm de caja + separación) sin comerse el cuerpo:
% el texto arranca a 4 cm del borde y el ancho de texto sigue siendo 16 cm.
\usepackage[a4paper, left=2.5cm, right=2.5cm, bottom=2.5cm,
            top=1.2cm, includehead, headheight=2.2cm, headsep=0.6cm]{geometry}

% --- Tipografía: Carlito, la misma decisión que en style.tex ---
% (clon métrico libre de Calibri; pdflatex no puede usar la real).
% No se replica el arreglo de versalitas de style.tex: aquí no hay
% bibliografía APA y ningún texto usa \textsc.
\usepackage[sfdefault]{carlito}

% --- Espaciado: interlineado sencillo + 6 pt entre párrafos (guía) ---
\usepackage[skip=6pt]{parskip}

% --- Tablas ---
% longtable: el registro y el plan de respuesta superan una página.
% AVISO ESPERADO: "ignored error: Infinite glue shrinkage" en cada salto de
% página de una longtable (bug conocido e inofensivo, latex3/latex2e#1907;
% ver el comentario largo en packages.tex de la memoria).
\usepackage{array}
\usepackage{longtable}
\usepackage{tabularx}
% \RaggedRight (ragged2e) en vez de \raggedright en columnas muy estrechas:
% el \raggedright de LaTeX desactiva de facto la partición silábica y una
% palabra larga («Identificado») desborda una columna de 1,2 cm; la
% variante de ragged2e mantiene el guionado.
\usepackage{ragged2e}
% Las columnas X se componen como m{...}: el contenido de cada celda queda
% centrado verticalmente (aspecto de formulario del diseño de referencia).
\renewcommand{\tabularxcolumn}[1]{m{#1}}
% Tipos de columna del diseño de referencia:
%   L/C = ancho fijo alineado izquierda/centro; Y/Z = X alineada izq./centro.
% En columnas estrechas se prefiere Y (raggedright) a X justificada: la
% justificación en 3-4 cm abre ríos de blancos.
\newcolumntype{L}[1]{>{\raggedright\arraybackslash}m{#1}}
\newcolumntype{C}[1]{>{\centering\arraybackslash}m{#1}}
\newcolumntype{Y}{>{\raggedright\arraybackslash}X}
\newcolumntype{Z}{>{\centering\arraybackslash}X}
% Filas con aire en todas las tablas (el diseño de referencia lo fijaba
% tabla a tabla; aquí una sola vez).
\renewcommand{\arraystretch}{1.25}

% --- Utilidades ---
\usepackage{xcolor}    % \markfield (resaltado) y \todonote (macros.tex)
\usepackage{amssymb}   % \square, \blacksquare, \circlearrowright
\usepackage{xspace}    % lo necesita \bukit (macros.tex)
\usepackage{enumitem}
\setlist[itemize]{leftmargin=*, topsep=2pt, itemsep=2pt}
\setlist[enumerate]{leftmargin=*, topsep=2pt, itemsep=2pt}

% --- Títulos según la guía de estilo, como en style.tex de la memoria ---
% (sin \chapter: la clase article no lo tiene; el resto, idéntico).
\usepackage{titlesec}
\titleformat{\section}
  {\normalfont\bfseries\fontsize{12pt}{14.4pt}\selectfont\raggedright}
  {\thesection.}{0.5em}{\MakeUppercase}
\titleformat{\subsection}
  {\normalfont\bfseries\fontsize{12pt}{14.4pt}\selectfont}
  {\thesubsection.}{0.5em}{}
\titleformat{\subsubsection}
  {\normalfont\itshape\fontsize{12pt}{14.4pt}\selectfont}
  {\thesubsubsection.}{0.5em}{}

% --- Paginación dinámica (regla estricta del autor, 26-07-2026) ---
% PROHIBIDO el salto de página manual también en estos documentos: el
% espacio lo garantizan estos ganchos (o un \Needspace* local ante un bloque
% indivisible), nunca un \newpage. Los .tex de referencia sí usaban
% \newpage; es la única desviación deliberada respecto a ellos.
\usepackage{needspace}
\AddToHook{cmd/section/before}{\Needspace*{5\baselineskip}}
\AddToHook{cmd/subsection/before}{\Needspace*{5\baselineskip}}
\AddToHook{cmd/subsubsection/before}{\Needspace*{4\baselineskip}}
\widowpenalty=10000
\clubpenalty=10000
\displaywidowpenalty=10000
\raggedbottom
% Estos formularios componen mucho texto en columnas de 2-4 cm: se concede
% a TeX un margen de estirado de emergencia antes de desbordar la línea
% (evita Overfull \hbox menores sin tocar el contenido).
\setlength{\emergencystretch}{1em}

% --- Datos globales y macros de la memoria (única fuente: no duplicar) ---
% =========================================================================
% Variables globales del documento.
% Cambiar un dato AQUÍ lo actualiza en todo el documento (portada,
% metadatos del PDF...). No escribir nunca estos datos "a mano" en otros
% archivos: usar siempre estos comandos.
% =========================================================================

% --- Trabajo ---
\newcommand{\thesisTitle}{BUKIT: red social orientada a objetivos}% TODO: título definitivo
\newcommand{\thesisSubtitle}{Memoria del Trabajo de Fin de Grado}%
\newcommand{\thesisDegree}{Grado en Ingeniería Informática de Gestión y Sistemas de Información}%
\newcommand{\thesisCourse}{2026--2027}% curso académico
\newcommand{\thesisDate}{Bilbao, \today}% TODO: <Lugar, día, mes, año> de la entrega/defensa
\newcommand{\thesisYear}{2027}% año de elaboración y defensa

% --- Personas (la plantilla del centro usa "Apellido1 Apellido2, Nombre") ---
\newcommand{\thesisAuthor}{Liviu Deleanu}% nombre natural (para metadatos del PDF)
\newcommand{\thesisAuthorFormal}{Deleanu, Liviu}% TODO: confirmar formato con dos apellidos
\newcommand{\thesisDirector}{Villamañe Gironés, Mikel}%
\newcommand{\thesisCoDirector}{Apellido1 Apellido2, Nombre}% TODO: solo si hay codirector/a

% --- Universidad ---
\newcommand{\thesisUniversity}{Universidad del País Vasco / Euskal Herriko Unibertsitatea}%
\newcommand{\thesisFaculty}{Escuela de Ingeniería de Bilbao}%

% --- Palabras clave del resumen trilingüe ---
\newcommand{\thesisKeywordsEs}{red social, objetivos, monorepo, TypeScript, PostgreSQL}% TODO
\newcommand{\thesisKeywordsEu}{sare soziala, helburuak, monorepo, TypeScript, PostgreSQL}% TODO: revisar euskera
\newcommand{\thesisKeywordsEn}{social network, goals, monorepo, TypeScript, PostgreSQL}% TODO

% =========================================================================
% Comandos propios. Definir aquí cualquier cosa que se repita en el texto:
% si mañana cambia cómo se escribe, se cambia en un único sitio.
% =========================================================================

% Nombre del producto, siempre con el mismo formato.
\newcommand{\bukit}{\textsc{Bukit}\xspace}

% Términos en inglés, en cursiva de forma consistente: \eng{backend}
\newcommand{\eng}[1]{\textit{#1}}

% Notas de trabajo, visibles en rojo en el PDF. Buscar "todonote" y
% eliminarlas todas antes de entregar.
\newcommand{\todonote}[1]{\textcolor{red}{\textbf{[TODO: #1]}}}


% --- Cabecera y pie (diseño de referencia, sin logos) ---
% La caja de cabecera replica la del .tex de referencia (celda ancha de
% contexto + celda estrecha de identidad). La celda derecha llevaba el logo
% del centro: aquí va texto porque los logos UPV/EHU siguen pendientes de
% la fase de diseño (ver style.tex); cuando existan, sustituir por
% \includegraphics. 11 + 4,9 = 15,9 cm < 16 cm de texto: sin desbordes.
\usepackage{fancyhdr}
\newcommand{\riesgosHeader}{%
  \setlength{\tabcolsep}{0pt}%
  \renewcommand{\arraystretch}{1}%
  \begin{tabular}{|p{11cm}|p{4.9cm}|}
    \hline
    \parbox[c][2cm][c]{11cm}{\centering\bfseries
      {\fontsize{16}{19}\selectfont TRABAJO DE FIN DE GRADO\\[4pt]Curso \thesisCourse}%
    } &
    \parbox[c][2cm][c]{4.9cm}{\centering
      {\fontsize{14}{17}\selectfont\bfseries\bukit}\\[3pt]
      {\small\thesisFaculty}%
    }\\
    \hline
  \end{tabular}%
}
\pagestyle{fancy}
\fancyhf{}
\renewcommand{\headrulewidth}{0pt}
\fancyhead[C]{\riesgosHeader}
\fancyfoot[C]{\thepage}
% La primera página usa estilo "plain": se redefine para que también lleve
% la cabecera (mismo truco que el .tex de referencia).
\fancypagestyle{plain}{%
  \fancyhf{}%
  \renewcommand{\headrulewidth}{0pt}%
  \fancyhead[C]{\riesgosHeader}%
  \fancyfoot[C]{\thepage}%
}

% --- Símbolos y campos de formulario del diseño de referencia ---
\newcommand{\fieldline}[1]{\underline{\hspace{#1}}}
\newcommand{\bigline}{\noindent\rule{\textwidth}{0.4pt}}
\newcommand{\checkbox}{\(\square\)}
\newcommand{\checkedbox}{\(\blacksquare\)}
% Símbolo ⟳ = "si se rellena este apartado, nueva versión del plan".
% El original lo resolvía con un SVG externo (paquete svg + Inkscape);
% aquí basta el símbolo de amssymb: cero dependencias nuevas.
\newcommand{\updatesymbol}{\raisebox{-0.2ex}{\(\circlearrowright\)}}
% Resaltado amarillo del identificador de riesgo en los títulos.
\newcommand{\markfield}[1]{\colorbox{yellow!35}{\strut #1}}

% --- Enlaces y metadatos del PDF (al final, como exige hyperref) ---
% El pdftitle lo fija cada documento con \hypersetup tras cargar esto.
\usepackage[hidelinks]{hyperref}
\hypersetup{pdfauthor={\thesisAuthor}}

\hypersetup{pdftitle={Seguimiento de riesgos — 16-08-2026 (S4)}}

\begin{document}

\begin{center}
  {\LARGE\bfseries Seguimiento de riesgos — S4}\\[6pt]
\end{center}

\textbf{Fecha:} 16-08-2026 · \textbf{Sprint:} S4 (10-08 \(\to\) 16-08-2026, cuarta retrospectiva) · \textbf{Responsable:} \thesisAuthor

Regla de revisión: \textbf{críticos} cada sprint; \textbf{significativos} cada 2 sprints; \textbf{moderados} cada 4 sprints; cualquier riesgo, \textbf{de inmediato}, si se observa su indicador clave.

\textbf{El símbolo \updatesymbol\ indica que, si se rellena ese apartado, hay que emitir una nueva versión completa del plan de riesgos} en \texttt{../plan-riesgos/} y registrar el cambio en la tabla de seguimiento del cap.~7 de la memoria.

Riesgos revisados — primera revisión completa del registro: \textbf{R1--R7 por cadencia} (críticos, semanal), \textbf{R8--R11 y R14 por cadencia} (significativos, cada dos sprints) y \textbf{R12--R13 por cadencia} (moderados, primera revisión, cada cuatro).
Datos del sprint: velocidad de S4 = 3 SP (US-05, cerrada el lunes 10-08) — exactamente la esperada para 6 horas declaradas con la equivalencia de 2,0 h/SP; horas reales \(\approx\)1~h~40 de 6 declaradas (28~\%), con el sprint entero concentrado en el lunes y \textbf{cinco días consecutivos sin actividad} (11--15/08) por causa sobrevenida (ver R10).

\bigline

\section{Medidas de mitigación (prevención)}

Para cada medida de los riesgos que tocan revisión en este seguimiento: ¿se está aplicando?
Si no, justificar y decidir si se mantiene/modifica/elimina.

% Bloques de 3-4 filas (tabularx no parte entre páginas; receta del S2).
\Needspace*{24\baselineskip}
\underline{\textbf{Riesgos críticos (R1--R4)}}

\begin{tabularx}{\textwidth}{|C{0.9cm}|Y|Y|C{2.2cm}|L{2.4cm}|}
\hline
\multicolumn{1}{|C{0.9cm}|}{\textbf{Id}} &
\multicolumn{1}{>{\centering\arraybackslash}X|}{\textbf{Medida de mitigación}} &
\multicolumn{1}{>{\centering\arraybackslash}X|}{\textbf{¿Aplicada?}} &
\multicolumn{1}{C{2.2cm}|}{\textbf{Justificación (del NO)}} &
\multicolumn{1}{C{2.4cm}|}{\textbf{Decisión}} \\
\hline
R1 & Toda idea nueva entra por el backlog. & \checkedbox\ Sí (R14 y sus mecanismos entraron por el registro de riesgos, no como alcance; US-05 siguió el orden).\newline \checkbox\ No & — & \checkedbox\ Mantener\newline \checkbox\ Modificar\newline \checkbox\ Eliminar \\
\hline
R2 & Medir velocidad y recalibrar en cada cierre. & \checkedbox\ Sí (cuarta medición: 3 SP; primera semana bajo la equivalencia recalibrada de 2,0, que la realidad confirma).\newline \checkbox\ No & — & \checkedbox\ Mantener\newline \checkbox\ Modificar\newline \checkbox\ Eliminar \\
\hline
R3 & Diario al cierre de cada sesión; memoria con hueco. & \checkedbox\ Sí (diario del lunes el mismo día; la memoria no tuvo cambios pendientes esta semana).\newline \checkbox\ No & — & \checkedbox\ Mantener\newline \checkbox\ Modificar\newline \checkbox\ Eliminar \\
\hline
R4 & Citar al escribir; pendientes con nota visible. & \checkedbox\ Sí (sin afirmaciones externas nuevas esta semana; ningún todonote nuevo sin registrar).\newline \checkbox\ No & — & \checkedbox\ Mantener\newline \checkbox\ Modificar\newline \checkbox\ Eliminar \\
\hline
\end{tabularx}

\Needspace*{24\baselineskip}
\underline{\textbf{Riesgos críticos (R5--R7) y significativo R8}}

\begin{tabularx}{\textwidth}{|C{0.9cm}|Y|Y|C{2.2cm}|L{2.4cm}|}
\hline
\multicolumn{1}{|C{0.9cm}|}{\textbf{Id}} &
\multicolumn{1}{>{\centering\arraybackslash}X|}{\textbf{Medida de mitigación}} &
\multicolumn{1}{>{\centering\arraybackslash}X|}{\textbf{¿Aplicada?}} &
\multicolumn{1}{C{2.2cm}|}{\textbf{Justificación (del NO)}} &
\multicolumn{1}{C{2.4cm}|}{\textbf{Decisión}} \\
\hline
R5 & Revisar todo antes de integrar; nada sin comprobaciones. & \checkedbox\ Sí (la revisión adversarial de US-05 cazó 4 defectos de seguridad reales antes del commit, con uno refutado con motivo).\newline \checkbox\ No & — & \checkedbox\ Mantener\newline \checkbox\ Modificar\newline \checkbox\ Eliminar \\
\hline
R6 & Investigación acotada ante tecnología nueva. & \checkedbox\ Sí (jsonwebtoken se eligió con comparación explícita frente a jose; sin descarrilamientos).\newline \checkbox\ No & — & \checkedbox\ Mantener\newline \checkbox\ Modificar\newline \checkbox\ Eliminar \\
\hline
R7 & Minimización desde el diseño (RGPD). & \checkedbox\ Sí (US-05 no añadió entidad ni dato personal alguno; los tokens no persisten).\newline \checkbox\ No & — & \checkedbox\ Mantener\newline \checkbox\ Modificar\newline \checkbox\ Eliminar \\
\hline
R8 & Horas por sesión en el diario; cada retro compara. & \checkedbox\ Sí (comparación hecha: 1 h 40 de 6; ver su ficha).\newline \checkbox\ No & — & \checkedbox\ Mantener\newline \checkbox\ Modificar\newline \checkbox\ Eliminar \\
\hline
\end{tabularx}

\Needspace*{24\baselineskip}
\underline{\textbf{Riesgos significativos (R9--R11 y R14)}}

\begin{tabularx}{\textwidth}{|C{0.9cm}|Y|Y|C{2.2cm}|L{2.4cm}|}
\hline
\multicolumn{1}{|C{0.9cm}|}{\textbf{Id}} &
\multicolumn{1}{>{\centering\arraybackslash}X|}{\textbf{Medida de mitigación}} &
\multicolumn{1}{>{\centering\arraybackslash}X|}{\textbf{¿Aplicada?}} &
\multicolumn{1}{C{2.2cm}|}{\textbf{Justificación (del NO)}} &
\multicolumn{1}{C{2.4cm}|}{\textbf{Decisión}} \\
\hline
R9 & Objetivos secundarios o versiones simuladas; capas gratuitas. & \checkedbox\ Sí (Render operativo sin incidencias; sin integraciones nuevas).\newline \checkbox\ No & — & \checkedbox\ Mantener\newline \checkbox\ Modificar\newline \checkbox\ Eliminar \\
\hline
R10 & Margen deliberado + documentación continua. & \checkedbox\ Sí (retomar hoy = leer el diario del lunes; el tema del sprint ya estaba entregado ANTES del parón: el margen ni se tocó).\newline \checkbox\ No & — & \checkedbox\ Mantener\newline \checkbox\ Modificar\newline \checkbox\ Eliminar \\
\hline
R11 & Todo en git con remoto y subidas frecuentes. & \checkedbox\ Sí (las dos tandas del lunes subidas el mismo día; árbol limpio toda la semana).\newline \checkbox\ No & — & \checkedbox\ Mantener\newline \checkbox\ Modificar\newline \checkbox\ Eliminar \\
\hline
R14 & Auditoría de dependencias en CI + Dependabot + casos de abuso en la DoD. & \checkedbox\ Sí y OBSERVABLE: el audit corrió en verde en cada push y PR; Dependabot generó dos tandas de PRs (09-08 y 16-08, con CI verde); US-05 cerró con 11 casos de abuso.\newline \checkbox\ No & — & \checkedbox\ Mantener\newline \checkbox\ Modificar\newline \checkbox\ Eliminar \\
\hline
\end{tabularx}

\Needspace*{16\baselineskip}
\underline{\textbf{Riesgos moderados (R12--R13, primera revisión por cadencia)}}

\begin{tabularx}{\textwidth}{|C{0.9cm}|Y|Y|C{2.2cm}|L{2.4cm}|}
\hline
\multicolumn{1}{|C{0.9cm}|}{\textbf{Id}} &
\multicolumn{1}{>{\centering\arraybackslash}X|}{\textbf{Medida de mitigación}} &
\multicolumn{1}{>{\centering\arraybackslash}X|}{\textbf{¿Aplicada?}} &
\multicolumn{1}{C{2.2cm}|}{\textbf{Justificación (del NO)}} &
\multicolumn{1}{C{2.4cm}|}{\textbf{Decisión}} \\
\hline
R12 & Capítulos a revisión conforme se cierran; plazo orientativo por envío. & \checkbox\ Sí\newline \checkedbox\ No: aún no se ha enviado nada al tutor. & La consulta pendiente (§3.5.2, anotada desde el 30-07) sigue sin enviarse; sin envío no hay ciclo que medir. & \checkedbox\ Mantener\newline \checkbox\ Modificar\newline \checkbox\ Eliminar \\
\hline
R13 & Cada solución como receta reproducible al resolverla. & \checkedbox\ Sí (recetas de la cadencia: graphify/Smart App Control en CLAUDE.md §3, migración no interactiva de Prisma, relanzamiento de Docker, diagnóstico en vivo de jobs de CI).\newline \checkbox\ No & — & \checkedbox\ Mantener\newline \checkbox\ Modificar\newline \checkbox\ Eliminar \\
\hline
\end{tabularx}

\Needspace*{6\baselineskip}
\underline{\textbf{Cambios realizados en las medidas de mitigación}} \updatesymbol

\begin{tabularx}{\textwidth}{|C{1.4cm}|Y|Y|Y|}
\hline
\multicolumn{1}{|C{1.4cm}|}{\textbf{Id}} &
\multicolumn{1}{>{\centering\arraybackslash}X|}{\textbf{Medida original}} &
\multicolumn{1}{>{\centering\arraybackslash}X|}{\textbf{Cambio realizado}} &
\multicolumn{1}{>{\centering\arraybackslash}X|}{\textbf{Justificación}} \\
\hline
— & — & — & — \\
\hline
\end{tabularx}

\Needspace*{6\baselineskip}
\underline{\textbf{Nuevas medidas de mitigación a añadir}} \updatesymbol

\begin{tabularx}{\textwidth}{|C{1.4cm}|Y|Y|}
\hline
\multicolumn{1}{|C{1.4cm}|}{\textbf{Id}} &
\multicolumn{1}{>{\centering\arraybackslash}X|}{\textbf{Nueva medida}} &
\multicolumn{1}{>{\centering\arraybackslash}X|}{\textbf{Justificación}} \\
\hline
— & — & — \\
\hline
\end{tabularx}

\bigline

\section{Riesgos a revisar en este seguimiento}

Para no alargar el documento sin señal, los riesgos \textbf{sin novedad y sin observación relevante} se agrupan en la tabla siguiente (indicador en su valor esperado, sin salto, estado «sin novedad»); los que dejaron señal esta semana (R8, R10, R12, R14) conservan su ficha completa a continuación.

\Needspace*{18\baselineskip}
\begin{tabularx}{\textwidth}{|C{0.9cm}|Y|C{2.6cm}|}
\hline
\multicolumn{1}{|C{0.9cm}|}{\textbf{Id}} &
\multicolumn{1}{>{\centering\arraybackslash}X|}{\textbf{Indicador: valor real (esperado 0 en todos)}} &
\multicolumn{1}{C{2.6cm}|}{\textbf{Nuevo estado}} \\
\hline
R1 & 0 — nada entró fuera del backlog; el sprint no ganó puntos. & Sin novedad \\
\hline
R2 & Indicador (3 primeros sprints) ya evaluado en S3; cuarta medición: 3 SP, la esperada con 2,0 h/SP — primera evidencia a favor de la recalibración. & Sin novedad \\
\hline
R3 & 0 — diario del lunes el mismo día; ningún capítulo pendiente de trabajo cerrado. & Sin novedad \\
\hline
R5 & 0 — decisiones de US-05 explicadas en el diario con alternativas; revisión adversarial con criterio (hallazgos aplicados Y refutados). & Sin novedad \\
\hline
R6 & 0 — sin historias al doble de tiempo; jsonwebtoken elegido con investigación acotada. & Sin novedad \\
\hline
R7 & 0 — US-05 sin datos personales nuevos (tokens sin persistir). & Sin novedad \\
\hline
R9 & 0 — sin cambios de precios/límites ni indisponibilidades; Render sin incidencias. & Sin novedad \\
\hline
R11 & 0 — todo subido el mismo lunes; sin anomalías del equipo. & Sin novedad \\
\hline
R13 & 0 — episodios de la cadencia (bloqueo de graphify por Smart App Control, caídas repetidas de Docker Desktop) resueltos en minutos y con receta; ninguno superó la hora del indicador. & Sin novedad \\
\hline
\end{tabularx}

% ------------------ R8 ------------------
\Needspace*{12\baselineskip}
\subsection*{Riesgo: \markfield{R8 — Capacidad real inferior a la declarada}}

\textbf{Indicadores (clave y de resolución)}

\begin{tabularx}{\textwidth}{|Z|Z|Z|Z|Z|}
\hline
\textbf{Indicador} & \textbf{Valor esperado} & \textbf{Valor real} & \textbf{¿Cumplido?} & \textbf{Acción / observaciones} \\
\hline
Horas <70~\% de las declaradas dos semanas consecutivas sin causa puntual identificada. & 0 & Primera semana bajo la capacidad RE-BASADA: 1 h 40 de 6 (28~\%), pero CON causa puntual identificada (la doble causa sobrevenida de R10). El episodio de S1--S3 quedó resuelto con la re-base y este es un cómputo nuevo contra las 6 horas. & \checkbox\ Sí\quad \checkedbox\ No & Vigilancia expresa para S5: si se repite el déficit contra las 6 horas, tocará valorar si la capacidad de verano admite más re-base o si el margen del calendario debe absorberlo. Matiz a favor: la velocidad del sprint fue la esperada (3 SP) — el plan re-basado aguantó una semana casi vacía. \\
\hline
\end{tabularx}

\textbf{Nueva probabilidad (1–3):} \fieldline{0.8cm} · \textbf{Nuevo impacto (1–3):} \fieldline{0.8cm} · \textbf{Nueva severidad:} \checkbox\ Crítico · \checkbox\ Significativo · \checkbox\ Moderado\ \updatesymbol

\textbf{¿Cerrar el riesgo?} \checkbox\ Sí · \checkedbox\ No\ \updatesymbol

\textbf{¿Se ha materializado?} \checkbox\ Sí · \checkedbox\ No\ \updatesymbol

\textbf{Nuevo estado} (cap.~7): \checkedbox\ Sin novedad · \checkbox\ Activo · \checkbox\ Materializado · \checkbox\ Resuelto · \checkbox\ Cerrado

\textbf{Plan de contingencia (solo si se ha materializado)}: no aplicable.

% ------------------ R10 ------------------
\Needspace*{12\baselineskip}
\subsection*{Riesgo: \markfield{R10 — Indisponibilidad prolongada del autor}}

\textbf{Indicadores (clave y de resolución)}

\begin{tabularx}{\textwidth}{|Z|Z|Z|Z|Z|}
\hline
\textbf{Indicador} & \textbf{Valor esperado} & \textbf{Valor real} & \textbf{¿Cumplido?} & \textbf{Acción / observaciones} \\
\hline
>3 días consecutivos sin poder trabajar por causa sobrevenida. & 0 & \textbf{5 días consecutivos} (11--15/08) por causa sobrevenida declarada por el autor: compromisos laborales más fatiga extrema por falta de sueño derivada de problemas personales. & \checkedbox\ Sí\quad \checkbox\ No & Segundo episodio (el primero, S1). Revisión dentro de esta retrospectiva; contingencia abajo. \\
\hline
\end{tabularx}

\textbf{Nueva probabilidad (1–3):} \underline{\ 1\ } · \textbf{Nuevo impacto (1–3):} \underline{\ 3\ } · \textbf{Nueva severidad:} \checkbox\ Crítico · \checkedbox\ Significativo · \checkbox\ Moderado\ \updatesymbol\ (sin cambio: dos episodios en siete sprints, ambos absorbidos; la evaluación se confirma)

\textbf{¿Cerrar el riesgo?} \checkbox\ Sí · \checkedbox\ No\ \updatesymbol

\textbf{¿Se ha materializado?} \checkedbox\ Sí · \checkbox\ No\ \updatesymbol

\textbf{Nuevo estado} (cap.~7): \checkbox\ Sin novedad · \checkbox\ Activo · \checkbox\ Materializado · \checkedbox\ Resuelto (materializado y superado en esta misma revisión, como en S1) · \checkbox\ Cerrado

\textbf{Plan de contingencia (solo si se ha materializado)}

\begin{tabularx}{\textwidth}{|Y|Y|Y|}
\hline
\multicolumn{1}{|>{\centering\arraybackslash}X|}{\textbf{Medida de contingencia}} &
\multicolumn{1}{>{\centering\arraybackslash}X|}{\textbf{¿Aplicada?}} &
\multicolumn{1}{>{\centering\arraybackslash}X|}{\textbf{Justificación (del NO)}} \\
\hline
Pausa formal + re-planificación del calendario (+ cambio de convocatoria en el extremo). & \checkedbox\ Sí, en su versión mínima (como en S1): pausa de facto sin re-planificación — el tema del sprint (US-05) se había entregado el lunes ANTES del parón, así que S4 cumplió su calendario y el margen no se consumió.\newline \checkbox\ No & — \\
\hline
\end{tabularx}

Reflexión sobre el funcionamiento de la contingencia:
el patrón «entregar la historia el lunes» convirtió un parón de cinco días en un sprint cumplido — la mecánica de flujo (historia pequeña, pronto en la semana) es en la práctica una mitigación extra de R10 que no estaba escrita.
Diferencia con S1: allí el parón dejó el sprint al límite (la historia se cerró el último día); aquí no costó nada.
Si los episodios se encadenan en S5--S6 (cierre de F1), tocará la conversación honesta sobre la capacidad real del verano.

% ------------------ R12 ------------------
\Needspace*{12\baselineskip}
\subsection*{Riesgo: \markfield{R12 — Ciclos de revisión con el tutor} (primera revisión por cadencia)}

\textbf{Indicadores (clave y de resolución)}

\begin{tabularx}{\textwidth}{|Z|Z|Z|Z|Z|}
\hline
\textbf{Indicador} & \textbf{Valor esperado} & \textbf{Valor real} & \textbf{¿Cumplido?} & \textbf{Acción / observaciones} \\
\hline
Consulta/capítulo >2 semanas sin respuesta, o revisión con cambios estructurales sobre capítulo cerrado. & 0 & 0 envíos realizados: no hay ciclo que medir. La consulta de §3.5.2 (anotada el 30-07) sigue SIN ENVIARSE — el riesgo no ha saltado, pero su mitigación (enviar pronto y por partes) tampoco se está ejercitando. & \checkbox\ Sí\quad \checkedbox\ No & Recomendación registrada: enviar la consulta de §3.5.2 esta semana; cada semana sin enviarla acerca el escenario del ciclo tardío que este riesgo describe. \\
\hline
\end{tabularx}

\textbf{Nueva probabilidad (1–3):} \fieldline{0.8cm} · \textbf{Nuevo impacto (1–3):} \fieldline{0.8cm} · \textbf{Nueva severidad:} \checkbox\ Crítico · \checkbox\ Significativo · \checkbox\ Moderado\ \updatesymbol

\textbf{¿Cerrar el riesgo?} \checkbox\ Sí · \checkedbox\ No\ \updatesymbol

\textbf{¿Se ha materializado?} \checkbox\ Sí · \checkedbox\ No\ \updatesymbol

\textbf{Nuevo estado} (cap.~7): \checkedbox\ Sin novedad · \checkbox\ Activo · \checkbox\ Materializado · \checkbox\ Resuelto · \checkbox\ Cerrado

\textbf{Plan de contingencia (solo si se ha materializado)}: no aplicable.

% ------------------ R14 ------------------
\Needspace*{12\baselineskip}
\subsection*{Riesgo: \markfield{R14 — Verificación de seguridad insuficiente} (primera revisión)}

\textbf{Indicadores (clave y de resolución)}

\begin{tabularx}{\textwidth}{|Z|Z|Z|Z|Z|}
\hline
\textbf{Indicador} & \textbf{Valor esperado} & \textbf{Valor real} & \textbf{¿Cumplido?} & \textbf{Acción / observaciones} \\
\hline
Historia con autenticación, autorización o datos personales cerrada sin casos de abuso, o vulnerabilidad alta sin resolver en el sprint. & 0 & 0 — la única historia del sprint (US-05, autenticación pura) cerró con 11 casos de abuso; la auditoría de dependencias pasó en verde en todos los pushes y PRs de la semana. & \checkbox\ Sí\quad \checkedbox\ No & La mitigación es ya observable: Dependabot generó dos tandas de PRs (5 abiertos, todos con la CI en verde), pendientes de triaje — los parche/menores pueden fusionarse; los mayores (globals 15\(\to\)17, @types/node 22\(\to\)26) piden revisión porque el proyecto fija Node por Volta. \\
\hline
\end{tabularx}

\textbf{Nueva probabilidad (1–3):} \fieldline{0.8cm} · \textbf{Nuevo impacto (1–3):} \fieldline{0.8cm} · \textbf{Nueva severidad:} \checkbox\ Crítico · \checkbox\ Significativo · \checkbox\ Moderado\ \updatesymbol

\textbf{¿Cerrar el riesgo?} \checkbox\ Sí · \checkedbox\ No\ \updatesymbol

\textbf{¿Se ha materializado?} \checkbox\ Sí · \checkedbox\ No\ \updatesymbol

\textbf{Nuevo estado} (cap.~7): \checkedbox\ Sin novedad · \checkbox\ Activo · \checkbox\ Materializado · \checkbox\ Resuelto · \checkbox\ Cerrado

\textbf{Plan de contingencia (solo si se ha materializado)}: no aplicable.

\bigline

\section{Reflexión desde el último seguimiento}

\textbf{¿Ha ocurrido algo reseñable desde el último seguimiento?} \checkedbox\ Sí · \checkbox\ No

Semana de dos caras: el lunes concentró TODO el sprint (R14 formalizado con su mitigación activada y US-05 entregada con sus 11 casos de abuso) y del martes al sábado no hubo actividad — cinco días que materializan R10 por segunda vez (causa sobrevenida: compromisos laborales más fatiga por problemas personales).
La lección buena queda escrita: con la historia entregada al inicio de la semana, el parón no costó calendario — velocidad 3 SP, la esperada, y F1 sigue en fecha (US-06 en S5, US-07 y cierre en S6).
La retrospectiva se celebró en su domingo, en sesión corta (40 minutos) y con lo prioritario cubierto.

\Needspace*{6\baselineskip}
\textbf{Nuevos riesgos identificados} \updatesymbol

\begin{tabularx}{\textwidth}{|C{1.4cm}|Y|C{0.8cm}|C{0.8cm}|C{1cm}|C{2.4cm}|}
\hline
\multicolumn{1}{|C{1.4cm}|}{\textbf{Id}} &
\multicolumn{1}{>{\centering\arraybackslash}X|}{\textbf{Descripción}} &
\multicolumn{1}{C{0.8cm}|}{\textbf{P}} &
\multicolumn{1}{C{0.8cm}|}{\textbf{I}} &
\multicolumn{1}{C{1cm}|}{\textbf{Exp.}} &
\multicolumn{1}{C{2.4cm}|}{\textbf{Severidad}} \\
\hline
— & — & — & — & — & — \\
\hline
\end{tabularx}

\textbf{¿Afecta algo a los objetivos o al alcance (cap.~3 de la memoria)?} \checkbox\ Sí · \checkedbox\ No

\textbf{Justificación:} el sprint cumplió su tema y su velocidad esperada; el episodio de R10 se absorbió sin tocar calendario ni alcance.

\textbf{Otros eventos relevantes:} primera cosecha de Dependabot (dos tandas de PRs con la CI en verde), pendiente de triaje con criterio anotado en R14; la línea de horas del diario del lunes fue confirmada por el autor (1 h).

\end{document}
